% Part: counterfactuals
% Chapter: introduction
% Section: material-conditional

\documentclass[../../../include/open-logic-section]{subfiles}

\begin{document}

\olfileid{cnt}{int}{mat}

\olsection{في الشرط المادي}

أبسط صيغ الشرط في الإنجليزية قول على مثال «إذا \dots، فإن \dots»،
ويحل في كل موضع من~\dots{} قول تام؛ ومنه «إذا كان الخادم هو الفاعل،
فإن البستاني بريء». ويُعلّم في مبادئ المنطق أن يُمثّل الشرط بالرابط
$\lif$: فإذا رمزنا إلى الجزأين اللذين ناب عنهما~\dots، مثلًا،
بـ!!{formula}~$!A$ و$!B$، كان رمز الشرط بجملته~$!A \lif !B$.

والرابط~$\lif$ \emph{دالة صدق}: فقيمة~$!A \lif !B$، وهي~$\True$
أو~$\False$، لا تتوقف إلا على قيمتي~$!A$ و$!B$. وذلك أن
$!A \lif !B$ يصدق إذا وفقط إذا كذب~$!A$ أو صدق~$!B$، ويكذب فيما
سوى ذلك. فإذا كان إسناد قيم الصدق~$\pAssign {\OLId{latin-ordinary}{v}}$، وضعنا
$\pSat{{\OLId{latin-ordinary}{v}}}{!A \lif !B}$ إذا وفقط إذا تحقق~$\pSat/{{\OLId{latin-ordinary}{v}}}{!A}$ أو
$\pSat{{\OLId{latin-ordinary}{v}}}{!B}$. والرابط~$\lif$ بهذه الدلالة هو المسمّى
\emph{الشرط المادي}.

وتلزم عن هذا الحدّ طائفة من الحقائق المنطقية الأولية. فأولها صحة مبرهنة
الاستنباط في الشرط المادي:
\begin{align}
  \text{إذا كان } {\OLId{greek-variable}{Gamma}}, !A \Entails !B \text{ فإن } {\OLId{greek-variable}{Gamma}} \Entails !A \lif !B
\end{align}
وهو، لكونه دالة صدق، يجعل~$!A \lif !B$ و$\lnot !A \lor !B$ متكافئتين:
\begin{align}
  !A \lif !B & \Entails \lnot !A \lor !B\\
  \lnot !A \lor !B & \Entails !A \lif !B
  \intertext{يستلزم تالي الشرط المادي الشرط نفسه، وكذلك يستلزمه نفي مقدمه:}
  !B & \Entails !A \lif !B\\
  \lnot !A & \Entails !A \lif !B
  \intertext{نفي الشرط المادي مكافئ لاقتران مقدمه بنفي تاليه: فإذا كان
    $!A \lif !B$ كاذبًا، كان~$!A \land \lnot !B$ صادقًا، والعكس:}
  \lnot(!A \lif !B) & \Entails !A \land \lnot !B\\
  !A \land \lnot !B & \Entails \lnot(!A \lif !B)
  \intertext{يدعم الشرط المادي قاعدة إثبات المقدّم:}
  !A, !A \lif !B & \Entails !B
  % تصحيح مصطلحي (COND-M1): سُمّيت قاعدة modus ponens إثبات المقدّم لا إثبات التالي.
  \intertext{يمكن تجميع الشروط المادية:}
  !A \lif !B,  !A \lif !C & \Entails !A \lif (!B \land !C)
  \intertext{يمكننا دائمًا تقوية المقدّم، أي إن الشرط \emph{رتيب}:}
  !A \lif !B & \Entails (!A \land !C) \lif !B
  \intertext{الشرط المادي متعدٍّ، أي إن قاعدة السلسلة صالحة:}
  !A \lif !B, !B \lif !C & \Entails !A \lif !C
  \intertext{الشرط المادي مكافئ لمقابله بالنقيض:}
  !A \lif !B & \Entails \lnot !B \lif \lnot !A\\
  \lnot !B \lif \lnot !A & \Entails !A \lif !B
\end{align}

وهذه وجوه استدلال نافعة لا حرج فيها في الرياضيات. إلا أن كتب الفلسفة
واللغة تورد كثيرًا من الأمثلة التي يُدّعى أنها تنقض نظائرها في الكلام
غير الرياضي. وهي تدل، إن سلمت، على أن الشرط المادي~$\lif$ ليس دائمًا
أنسب الروابط لتمثيل العبارة الإنجليزية التي صورتها «إذا \dots، فإن
\dots».

\end{document}

